% ============================================================================
% BAB V - KESIMPULAN DAN KENDALA
% ============================================================================
\chapter{KESIMPULAN DAN KENDALA}

\section{Kesimpulan}

Berdasarkan implementasi dan pengujian yang telah dilakukan, dapat
disimpulkan beberapa hal sebagai berikut:

\subsection{Keberhasilan Implementasi Replikasi Database}

Replikasi database PostgreSQL berhasil diimplementasikan menggunakan
metode \textit{streaming replication} asinkron antara 1 (satu)
PostgreSQL Master dan 2 (dua) PostgreSQL Replica. Konfigurasi meliputi
pengaturan \texttt{wal\_level = replica},
\texttt{max\_wal\_senders}, \texttt{hot\_standby}, serta pembuatan
user replikasi \texttt{repluser} dengan hak akses
\texttt{REPLICATION}. Inisialisasi Replica dilakukan secara otomatis
melalui script \texttt{entrypoint.sh} yang menjalankan
\texttt{pg\_basebackup} dengan flag \texttt{-R} untuk membuat
\texttt{standby.signal} dan \texttt{primary\_conninfo}. Seluruh
proses replikasi berjalan tanpa memerlukan intervensi manual setelah
konfigurasi awal diterapkan. Kedua Replica berhasil melakukan
sinkronisasi data dari Master secara kontinu melalui koneksi
\textit{WAL sender} yang persisten. Dengan adanya dua Replica, sistem
memiliki redundansi data yang memadai untuk mendukung mekanisme
\textit{failover} ketika salah satu \textit{backend} mengalami
gangguan.

\subsection{Keberhasilan Mekanisme Read-Write Splitting}

Mekanisme \textit{read-write splitting} berhasil diimplementasikan
melalui Pgpool-II dengan konfigurasi
\texttt{master\_slave\_mode = on},
\texttt{load\_balance\_mode = on}, dan
\texttt{master\_slave\_sub\_mode = 'stream'}. Hasil pengujian
membuktikan bahwa seluruh query \textit{write} (\texttt{INSERT},
\texttt{UPDATE}, \texttt{DELETE}) diproses oleh PostgreSQL Master,
sementara seluruh query \textit{read} (\texttt{SELECT})
didistribusikan ke PostgreSQL Replica. \textit{Load balancing}
berjalan dengan baik, di mana masing-masing Replica mendapat porsi query
\texttt{SELECT} yang seimbang sesuai bobot yang ditentukan
(\texttt{backend\_weight = 1}). Master dengan bobot \texttt{0}
sama sekali tidak menerima query \texttt{SELECT}, memastikan beban
operasi \textit{read} tidak mengganggu operasi \textit{write}. Bukti
log dari Pgpool-II (\texttt{log\_per\_node\_statement}) menunjukkan
setiap statement SQL tercatat pada node yang tepat sesuai
peruntukannya. Hal ini mengonfirmasi bahwa mekanisme \textit{routing}
otomatis Pgpool-II berfungsi dengan benar tanpa kesalahan
pengalihan query.

\subsection{Manfaat Penggunaan Pgpool-II}

Penggunaan Pgpool-II sebagai \textit{middleware} memberikan beberapa
manfaat signifikan. Pertama, \textbf{\textit{Connection Pooling}}
memungkinkan Pgpool-II mengelola \textit{pool} koneksi ke
\textit{backend} PostgreSQL sehingga mengurangi \textit{overhead}
pembukaan dan penutupan koneksi berulang yang dapat membebani server
database. Kedua, \textbf{\textit{Transparent Read-Write Splitting}}
membuat aplikasi tidak perlu memodifikasi kode untuk memisahkan query
\textit{read} dan \textit{write} karena Pgpool-II secara otomatis
mengarahkan query berdasarkan jenis operasi SQL. Ketiga,
\textbf{\textit{Load Balancing}} mendistribusikan query
\texttt{SELECT} ke beberapa Replica sehingga mencegah satu Replica
menjadi \textit{bottleneck} ketika traffic tinggi. Keempat,
\textbf{\textit{High Availability}} dijamin melalui mekanisme
\textit{health check} dan \textit{failover} yang memastikan aplikasi
tetap berjalan meskipun salah satu \textit{backend} mengalami
kegagalan. Kelima, \textbf{\textit{Single Entry Point}}
menyederhanakan konfigurasi koneksi karena aplikasi hanya perlu
terhubung ke satu alamat (Pgpool-II) tanpa perlu mengetahui topologi
database di belakangnya.

\subsection{Peningkatan Efisiensi Distribusi Query Database}

Dengan arsitektur \textit{read-write splitting}, beban query database
terdistribusi secara signifikan. Master hanya menangani operasi
\textit{write} (\texttt{INSERT}, \texttt{UPDATE}, \texttt{DELETE}),
sehingga memiliki lebih banyak sumber daya untuk menjamin integritas
data dan kecepatan komit transaksi. Sementara itu, Replica 1 dan
Replica 2 menangani seluruh operasi \textit{read} (\texttt{SELECT}),
yang merupakan mayoritas beban pada aplikasi \textit{e-commerce}
seperti Komposin. Distribusi ini secara signifikan meningkatkan
performa aplikasi karena operasi \textit{read} yang berat (seperti
laporan transaksi, pencarian produk, dan tampilan katalog) tidak lagi
membebani Master yang harus memproses operasi \textit{write} yang
kritis. Pengujian juga menunjukkan bahwa penambahan Replica dapat
dilakukan secara horizontal untuk meningkatkan kapasitas \textit{read}
lebih lanjut tanpa memengaruhi konfigurasi Master.

\section{Kendala dan Solusi}

Selama proses implementasi, beberapa kendala teknis ditemukan dan
diselesaikan. Berikut adalah kendala yang umum terjadi pada
implementasi arsitektur sejenis beserta analisis penyebab dan
solusinya:

\subsection{Kendala yang Mungkin Ditemukan}

\begin{enumerate}
    \item \textbf{Koneksi Antar Container Gagal.} \textit{Container}
    Pgpool-II tidak dapat terhubung ke PostgreSQL Master karena
    perbedaan jaringan atau kesalahan \textit{hostname resolution}.
    \begin{itemize}
        \item \textbf{Penyebab:} \textit{Container} belum tergabung
        dalam \textit{network} yang sama, atau \texttt{depends\_on}
        belum memastikan Master siap menerima koneksi.
        \item \textbf{Solusi:} Memastikan semua \textit{container}
        menggunakan \texttt{bdt-network} yang sama, menambahkan
        \texttt{healthcheck} pada Master, dan menggunakan
        \texttt{depends\_on} dengan
        \texttt{condition: service\_healthy}. Penggunaan IP statis
        juga membantu menghindari masalah resolusi hostname.
    \end{itemize}

    \item \textbf{Replika Gagal Melakukan pg\_basebackup.}
    \begin{itemize}
        \item \textbf{Penyebab:} Master belum siap, user replikasi
        belum dibuat, atau autentikasi ditolak oleh
        \texttt{pg\_hba.conf}.
        \item \textbf{Solusi:} Memastikan \texttt{init.sql} dijalankan
        sebelum Replica mencoba koneksi, serta mengonfigurasi
        \texttt{pg\_hba.conf} dengan baris
        \texttt{host replication repluser 172.25.0.0/24 md5}. Script
        \texttt{entrypoint.sh} juga perlu memeriksa keberadaan file
        \texttt{replica\_initialized} untuk menghindari pengulangan
        \texttt{pg\_basebackup} yang tidak perlu.
    \end{itemize}

    \item \textbf{Replikasi Tertunda (Replication Lag).}
    \begin{itemize}
        \item \textbf{Penyebab:} Beban \textit{write} yang tinggi atau
        koneksi jaringan antar \textit{container} lambat.
        \item \textbf{Solusi:} Meningkatkan \texttt{wal\_keep\_size}
        pada Master dan \texttt{hot\_standby\_feedback = on} pada
        Replica. Untuk kasus parah, mempertimbangkan
        \textit{synchronous replication} dengan mengaktifkan
        \texttt{synchronous\_standby\_names}.
    \end{itemize}

    \item \textbf{Port Database Terblokir Firewall Host.}
    \begin{itemize}
        \item \textbf{Penyebab:} Firewall \textit{host} (UFW/iptables)
        memblokir port yang digunakan \textit{container}
        (5433--5436).
        \item \textbf{Solusi:} Menambahkan aturan \texttt{ufw allow}
        untuk port yang diperlukan, atau menggunakan \texttt{expose}
        saja tanpa \texttt{ports} di \texttt{docker-compose.yml} jika
        akses hanya dibutuhkan antar \textit{container}.
    \end{itemize}

    \item \textbf{Kesalahan Konfigurasi pada pgpool.conf.}
    \begin{itemize}
        \item \textbf{Penyebab:} Kesalahan sintaks atau parameter
        yang tidak sesuai versi Pgpool-II.
        \item \textbf{Solusi:} Memvalidasi konfigurasi dengan membaca
        log Pgpool-II (\texttt{docker compose logs pgpool}) dan
        memastikan kompatibilitas parameter dengan versi image yang
        digunakan. Environment variable di \texttt{docker-compose.yml}
        juga perlu diperiksa karena di-\textit{override} ke
        konfigurasi.
    \end{itemize}

    \item \textbf{Permission Denied pada File Konfigurasi.}
    \begin{itemize}
        \item \textbf{Penyebab:} File konfigurasi yang di-\textit{mount}
        ke \textit{container} memiliki \textit{permission} yang tidak
        sesuai.
        \item \textbf{Solusi:} Mengatur \textit{permission} file
        konfigurasi (misalnya \texttt{chmod 644}) atau menggunakan
        \texttt{:ro} (\textit{read-only}) pada \textit{volume mount}
        untuk mencegah modifikasi oleh \textit{container}.
    \end{itemize}

    \item \textbf{SELECT Query Tetap ke Master.}
    \begin{itemize}
        \item \textbf{Penyebab:} \texttt{backend\_weight0} tidak diset
        ke \texttt{0} atau \texttt{load\_balance\_mode} tidak aktif.
        \item \textbf{Solusi:} Memastikan
        \texttt{PGPOOL\_PARAMS\_BACKEND\_WEIGHT0=0} dan
        \texttt{PGPOOL\_PARAMS\_LOAD\_BALANCE\_MODE=on} di environment
        Pgpool-II. Verifikasi dengan \texttt{SHOW pool\_nodes} untuk
        melihat \texttt{select\_cnt} pada Master.
    \end{itemize}

    \item \textbf{Aplikasi Tidak Dapat Terhubung ke Database.}
    \begin{itemize}
        \item \textbf{Penyebab:} \texttt{DB\_HOST} mengarah langsung
        ke Master, bukan ke Pgpool-II, atau konfigurasi
        \texttt{pool\_hba.conf} menolak koneksi.
        \item \textbf{Solusi:} Memastikan \texttt{DB\_HOST=pgpool} di
        \texttt{.env} dan \texttt{pool\_hba.conf} mengizinkan koneksi
        dari jaringan aplikasi. Koneksi dapat diuji dengan
        \texttt{docker exec bdt-app php artisan tinker} dan
        menjalankan query sederhana.
    \end{itemize}

    \item \textbf{Permission Denied pada Folder Unggah Gambar (/var/www/public/img).}
    \begin{itemize}
        \item \textbf{Penyebab:} Kontainer aplikasi Laravel (\texttt{bdt-app})
        berjalan dengan user \texttt{www-data} untuk proses web server Nginx
        dan PHP-FPM. Direktori \texttt{public/img} hasil \textit{copy} saat build
        dimiliki oleh user \texttt{root}, sehingga PHP-FPM tidak memiliki izin
        menulis data dan memicu error \textit{Unable to write in the "/var/www/public/img" directory}.
        \item \textbf{Solusi:} Mengubah kepemilikan folder secara rekursif menjadi
        \texttt{www-data:www-data} menggunakan perintah \texttt{chown -R} dan
        mengatur permission menjadi \texttt{775}. Perubahan ini kemudian ditambahkan
        ke dalam file \texttt{Dockerfile} agar terkonfigurasi secara otomatis saat build ulang.
    \end{itemize}
\end{enumerate}

\subsection{Proses Analisis dan Penyelesaian Masalah}

Pendekatan sistematis yang digunakan dalam menganalisis dan
menyelesaikan kendala meliputi lima langkah utama. Pertama,
\textbf{Analisis Log} dilakukan terhadap setiap layanan melalui
\texttt{docker compose logs [service]} untuk mengidentifikasi akar
permasalahan. Kedua, \textbf{Isolasi Masalah} dengan memeriksa apakah
masalah berasal dari jaringan (uji \texttt{ping}), konfigurasi
(validasi file \texttt{.conf}), autentikasi (periksa
\texttt{pg\_hba.conf}), atau data. Ketiga, \textbf{Pengujian
Bertahap} dilakukan dari layer terbawah (koneksi Master $\rightarrow$
Replica) hingga teratas (Aplikasi $\rightarrow$ Pgpool-II) untuk
mempersempit cakupan debugging. Keempat, \textbf{Verifikasi Manual}
menggunakan \texttt{docker exec} untuk mengakses \textit{container}
dan menjalankan perintah diagnostik seperti \texttt{ping},
\texttt{psql}, dan \texttt{SHOW pool\_nodes}. Kelima,
\textbf{Pencatatan Solusi} memastikan setiap solusi yang berhasil
diterapkan didokumentasikan untuk referensi di masa mendatang dan
mempercepat proses troubleshooting jika masalah serupa muncul kembali.
